\documentclass[14pt,dvipdfmx]{jsarticle}
\input{lecture-preamble.tex}

\begin{document}

\lectureheader[5]{数学Ⅱ}{6}{微分法と積分法}{1}{微分係数と導関数}{2-B}{関数の微分}
  {186,187}{定数倍と和・差の性質を使って, 整式を微分しよう}

% ============ 1ページ目（表・左） ============
\vspace{-2.0zw}%
\heading{〜前時の復習〜}
\vspace{-0.5em}
\begin{reviewbox}{微分の公式}
  $k$ を定数とする.
  \par\vspace{0.15em}
  {\renewcommand{\fitblankpad}{0.5zw}%
  $(x^{n})'=\fitblankbf{nx^{n-1}}$,\qquad
  $(c)'=\fitblankbf{0}$,\qquad
  $(kf(x))'=\fitblankbf{kf'(x)}$}
  \par\vspace{0.25em}
  {\renewcommand{\fitblankpad}{0.5zw}%
  $(f(x)+g(x))'=\fitblankbf{f'(x)+g'(x)}$,\qquad
  $(f(x)-g(x))'=\fitblankbf{f'(x)-g'(x)}$}
\end{reviewbox}

\vspace{0.1em}
\begin{exampleprob}{1}
  次の関数を微分せよ.
  \hspace*{2zw}$y=x(x+1)(x-2)$
\end{exampleprob}
\vspace{0.5em}
\begin{answerbox}{}[9]
  $x(x+1)(x-2)=x(x^{2}-x-2)=x^{3}-x^{2}-2x$\\[0.3em]
  よって $y=x^{3}-x^{2}-2x$\\[0.25em]
  したがって $y'=3x^{2}-2x-2$
  \boxanswer{$3x^{2}-2x-2$}
\end{answerbox}

\vspace{0.05em}
\begin{practiceprob}{9}
  次の関数を微分せよ.
  \begin{tasks}[label=(\arabic*), label-width=1.6zw, item-indent=2zw, after-item-skip=0.08em](2)
    \task $y=(x+2)(x+3)$
    \task $y=3(x-2)^{2}$
  \end{tasks}
\end{practiceprob}
\vspace{-0.45em}
\noindent
\begin{minipage}[t]{0.49\linewidth}
\begin{answerbox}{(1)}[10]
  $y=(x+2)(x+3)$\\[0.15em]
  $=x^{2}+5x+6$\\[0.15em]
  $y'=2x+5$
\end{answerbox}
\end{minipage}\hfill
\begin{minipage}[t]{0.49\linewidth}
\begin{answerbox}{(2)}[10]
  $y=3(x^{2}-4x+4)$\\[0.15em]
  $=3x^{2}-12x+12$\\[0.15em]
  $y'=6x-12$
\end{answerbox}
\end{minipage}

\vfill
\nextpage

% ============ 2ページ目（表・右） ============
\vspace{-1.4zw}%
\noindent{\bfseries 練習\,9}\quad（続き）
\begin{tasks}[label=(\arabic*), label-width=1.6zw, item-indent=2zw, after-item-skip=0.08em, start=3](2)
  \task $y=x(x+2)(x-2)$
  \task $y=2x(x+1)(x-3)$
\end{tasks}
\vspace{-0.2em}
\noindent
\begin{minipage}[t]{0.49\linewidth}
\begin{answerbox}{(3)}[10]
  $y=x(x^{2}-4)$\\[0.15em]
  $=x^{3}-4x$\\[0.15em]
  $y'=3x^{2}-4$
\end{answerbox}
\end{minipage}\hfill
\begin{minipage}[t]{0.49\linewidth}
\begin{answerbox}{(4)}[10]
  $y=2x(x^{2}-2x-3)$\\[0.1em]
  $=2x^{3}-4x^{2}-6x$\\[0.15em]
  $y'=6x^{2}-8x-6$
\end{answerbox}
\end{minipage}

\vspace{0.15em}
\begin{exampleprob}{2}
  関数 $f(x)=-x^{2}+4x+1$ の導関数を利用して,
  $x=-1$ における $f(x)$ の微分係数を求めよ.
\end{exampleprob}
\vspace{0.5em}
\begin{answerbox}{}[8.5]
  $f(x)$ を微分すると $f'(x)=-2x+4$\\[0.3em]
  求める微分係数は $f'(-1)=-2\cdot(-1)+4=6$
  \boxanswer{$6$}
\end{answerbox}

\vspace{0.5em}
\begin{practiceprob}{10}
  関数 $f(x)=x^{3}-3x^{2}+3$ について,
  次の $x$ の値における微分係数を求めよ.
  \begin{tasks}[label=(\arabic*), label-width=1.6zw, item-indent=2zw, after-item-skip=0.02em](3)
    \task $x=2$
    \task $x=0$
    \task $x=-2$
  \end{tasks}
\end{practiceprob}
\vspace{1em}
\noindent
$f'(x)=\fitblankbf{3x^{2}-6x}$ であるから
\par\vspace{0.5em}
\begin{answerbox}{(1)}[4.5]
  $f'(2)=3\cdot 4-6\cdot 2=0$
\end{answerbox}
\vspace{-0.25em}
\begin{answerbox}{(2)}[4.5]
  $f'(0)=0$
\end{answerbox}
\vspace{-0.25em}
\begin{answerbox}{(3)}[4.5]
  $f'(-2)=3\cdot 4-6\cdot(-2)=24$
\end{answerbox}

\vfill
\nextpage

% ============ 3ページ目（裏・左） ============
\vspace{-1.4zw}%
\heading{いろいろな関数の導関数}
\vspace{-0.4em}
変数が $x$, $y$ 以外の文字で表される関数についても, 同様に導関数を考える.
たとえば, $t$ の関数 $s=f(t)$ の導関数は
\par\vspace{0.2em}
\sbox{\fitblanksavebox}{\answerstyle{$\displaystyle\dfrac{d}{dt}f(t)$}}%
\newcommand{\notationblank}[1]{%
  \underline{\makebox[\dimexpr\wd\fitblanksavebox+2zw\relax][c]{%
    \ifshowanswer\answerstyle{$\displaystyle #1$}%
    \else\phantom{\answerstyle{$\displaystyle #1$}}\fi}}}%
\notationblank{s'},\quad \notationblank{f'(t)},\quad
\notationblank{\dfrac{ds}{dt}},\quad \notationblank{\dfrac{d}{dt}f(t)}
などで表すこともできる.

\vspace{1em}
\noindent{\bfseries 例\,7}\quad
$g$ は定数とする. $s=\dfrac{1}{2}gt^{2}$ を $t$ の関数とみて $t$ で微分すると
\par\vspace{0.15em}
\hspace*{2zw}$s'=\dfrac{1}{2}g\cdot 2t=\fitblankbf{gt}$

\vspace{1em}
{\settowidth{\dimen0}{\textbf{補足}\quad}%
\noindent\hangindent=\dimen0\hangafter=1
\textbf{補足}\quad
$s$ が $t$ の関数であり, その導関数を求めていることを式で明確にしたいときは,
$s'$ の代わりに $\dfrac{ds}{dt}$ を用いることがある.\par}

\vspace{1em}
\begin{practiceprob}{11}
  次の $t$ の関数を微分せよ. ただし, $a$, $b$ は定数とする.
  \begin{tasks}[label=(\arabic*), label-width=1.6zw, item-indent=2zw, after-item-skip=0.06em](2)
    \task $s=3t^{2}-4t+2$
    \task $f(t)=at^{2}+bt^{2}$
  \end{tasks}
\end{practiceprob}
\vspace{-0.55em}
\noindent
\begin{minipage}[t]{0.49\linewidth}
\begin{answerbox}{(1)}[4.8]
  $s'=6t-4$
\end{answerbox}
\end{minipage}\hfill
\begin{minipage}[t]{0.49\linewidth}
\begin{answerbox}{(2)}[4.8]
  $f'(t)=2at+2bt$
\end{answerbox}
\end{minipage}

\vspace{0.5em}
\begin{practiceprob}{12}
  半径 $r$ の球の体積を $V$, 表面積を $S$ とすると,
  $V=\dfrac{4}{3}\pi r^{3}$, $S=4\pi r^{2}$ である.
  $V$ と $S$ を $r$ の関数とみて, それぞれ $r$ で微分せよ.
\end{practiceprob}
\vspace{0.5em}
\begin{answerbox}{$V$}[5]
  $\dfrac{dV}{dr}=\dfrac{4}{3}\pi\cdot 3r^{2}=4\pi r^{2}$
\end{answerbox}
\vspace{-0.25em}
\begin{answerbox}{$S$}[5]
  $\dfrac{dS}{dr}=8\pi r$
\end{answerbox}

\vfill
\nextpage

% ============ 4ページ目（裏・右） ============
\vspace{-1.4zw}%
% \heading{研究\quad 関数 $x^{n}$ の導関数}
% \vspace{-0.45em}
% 183ページで学んだ次の導関数の公式を証明してみよう.
% \par\vspace{0.25em}
% \hspace*{2zw}$n$ を正の整数とするとき\quad
% $(x^{n})'=\fitblankbf{nx^{n-1}}$

% \vspace{0.35em}
% \textbf{【証明】}\quad
% 導関数の定義より
% \[
%   (x^{n})'=\lim_{h \to 0}\frac{(x+h)^{n}-x^{n}}{h}
% \]
% 二項定理により
% \begin{align*}
%   (x+h)^{n}
%   &={}_{n}\mathrm{C}_{0}x^{n}+{}_{n}\mathrm{C}_{1}x^{n-1}h
%     +{}_{n}\mathrm{C}_{2}x^{n-2}h^{2}+\cdots+{}_{n}\mathrm{C}_{n}h^{n}\\
%   &=x^{n}+nx^{n-1}h+\frac{1}{2}n(n-1)x^{n-2}h^{2}+\cdots+h^{n}
% \end{align*}
% よって
% \[
%   (x+h)^{n}-x^{n}
%   =nx^{n-1}h+\frac{1}{2}n(n-1)x^{n-2}h^{2}+\cdots+h^{n}
% \]
% 両辺を $0$ でない数 $h$ で割ると
% \begin{align*}
%   \frac{(x+h)^{n}-x^{n}}{h}
%   &=nx^{n-1}+\frac{1}{2}n(n-1)x^{n-2}h+\cdots+h^{n-1}
% \end{align*}
% したがって
% \[
%   \lim_{h \to 0}\frac{(x+h)^{n}-x^{n}}{h}
%   =\fitblankbf{nx^{n-1}}
% \]
% すなわち $(x^{n})'=\fitblankbf{nx^{n-1}}$

\noindent{\bfseries 演習問題}\quad
次の関数 $f(x)$ を微分せよ.
\par\vspace{-0.15em}
\begin{tasks}[label=(\arabic*), label-width=1.6zw, item-indent=2zw, after-item-skip=0.12em](2)
  \task $f(x)=2x^{2}+3x+4$
  \task $f(x)=-3x^{3}+5x^{2}-2x+1$
  \task $f(x)=(x-2)(3x-4)$
  \task $f(x)=(3x+1)^{3}$
\end{tasks}
\vspace{0.1em}
\noindent
\begin{minipage}[t]{0.49\linewidth}
\begin{answerbox}{(1)}[8]
  $f'(x)=2\cdot 2x+3\cdot 1=4x+3$
  \boxanswer{$4x+3$}
\end{answerbox}
\end{minipage}\hfill
\begin{minipage}[t]{0.49\linewidth}
\begin{answerbox}{(2)}[8]
  $f'(x)=-3\cdot 3x^{2}+5\cdot 2x-2$\\[0.15em]
  $=-9x^{2}+10x-2$
  \boxanswer{$-9x^{2}+10x-2$}
\end{answerbox}
\end{minipage}

\noindent
\begin{minipage}[t]{0.49\linewidth}
\begin{answerbox}{(3)}[11]
  $f(x)=3x^{2}-10x+8$\\[0.15em]
  $f'(x)=6x-10$
  \boxanswer{$6x-10$}
\end{answerbox}
\end{minipage}\hfill
\begin{minipage}[t]{0.49\linewidth}
\begin{answerbox}{(4)}[11]
  $f(x)=27x^{3}+27x^{2}+9x+1$\\[0.15em]
  $f'(x)=81x^{2}+54x+9$
  \boxanswer{$81x^{2}+54x+9$}
\end{answerbox}
\end{minipage}

\vfill
\begin{summarybox}
  \begin{enumerate}[label=\textbf{\arabic*}, leftmargin=1.5zw, itemsep=0.18em, topsep=0.1em]
    \item 積の形は\fitblankbf{\text{展開して}}から微分する.
    \item 微分係数 $f'(a)$ は, $f'(x)$ に
          $x=\fitblankbf{a}$ を代入して求める.
    \item $t$ の関数でも同様に微分でき,
          $s'$, $\dfrac{ds}{dt}$ などで表す.
  \end{enumerate}
\end{summarybox}

\vspace{0.15em}
\nexttime{接線の方程式を, 微分を使って求めてみよう!}

\end{document}
